\chapter{Wer in den Daten ist}

\section{Die Frage}

\begin{merksatz}
\boxtitel{Löhne gibt es nur für Erwerbstätige}
Wie stark erhöht ein Jahr Bildung den Lohn? Man regressiert den Lohn auf die Bildungsjahre --- aber Löhne beobachtet man nur bei denen, die arbeiten.

\begin{center}
{\small
\renewcommand{\arraystretch}{1.48}
\begin{tabular}{@{}P{3.6cm}P{2.6cm}@{}}
\toprule
\textbf{Verfahren} & \textbf{je Bildungsjahr}\\
\midrule
KQ, nur Erwerbstätige & $0{,}048$\\
\midrule
wahrer Wert & $\mathbf{0{,}080}$\\
\bottomrule
\end{tabular}}
\end{center}

\begin{center}
\begin{minipage}{0.9\textwidth}
\textbf{Die Regression unterschätzt um vierzig Prozent --- obwohl die Bildung in der Gesamtbevölkerung exogen ist.}
\end{minipage}
\end{center}

\begin{center}
\begin{minipage}{0.88\textwidth}
\kursiv{\textbf{Wer wenig Bildung hat und trotzdem arbeitet, hat meist einen besonderen Grund} --- eine hohe unbeobachtete Produktivität, einen Partner ohne Einkommen, keine Kinder. Unter den Erwerbstätigen stehen den Geringgebildeten also überdurchschnittlich gute Löhne gegenüber, und die Steigung wird flacher. Das ist die Auswahl aus Satz 5.4, nun nach einer eigenen Regel.}
\end{minipage}
\end{center}
\end{merksatz}

\begin{definition}
\boxtitelm{Definition 6.1 --- das Auswahlmodell}
Zwei Gleichungen mit gemeinsam normalverteilten Störungen:
\[ d^{*} = z^{\top}\gamma + \varepsilon_1 \qquad \text{(Teilnahme: } d = \mathbf 1\{d^{*}>0\}\text{)} \]
\[ y = x^{\top}\beta + \varepsilon_2 \qquad \text{(Ergebnis, beobachtet nur bei } d = 1\text{)} \]
mit $\operatorname{Corr}(\varepsilon_1,\varepsilon_2) = \rho$ und $\operatorname{Var}(\varepsilon_2) = \sigma^{2}$.
\end{definition}

\begin{merksatz}
\boxtitel{Satz 6.2 --- Heckmans Korrektur}
Unter den Teilnehmern gilt
\[ \mathbb{E}[y\mid x, d=1] = x^{\top}\beta + \rho\,\sigma\,\lambda(z^{\top}\gamma), \qquad \lambda(t) = \frac{\varphi(t)}{\Phi(t)} . \]

\schritt{Beweis}
Bedingt auf $\varepsilon_1 > -z^{\top}\gamma$ hat $\varepsilon_2$ wegen der gemeinsamen Normalverteilung den Erwartungswert $\rho\sigma\,\mathbb{E}[\varepsilon_1\mid\varepsilon_1>-z^{\top}\gamma]$, und der abgeschnittene Erwartungswert ist $\lambda(z^{\top}\gamma)$ --- wie in Satz 5.4. $\blacksquare$

\begin{center}
\begin{minipage}{0.9\textwidth}
\textbf{Die Auswahl wirkt wie eine ausgelassene Größe --- und diese Größe lässt sich berechnen.}
\end{minipage}
\end{center}

\schritt{Das zweistufige Verfahren}
\begin{center}
{\small
\renewcommand{\arraystretch}{1.5}
\begin{tabular}{@{}P{2.2cm}P{6.4cm}@{}}
\toprule
\textbf{Stufe 1} & Probit der Teilnahme auf $z$; daraus $\hat\lambda_i$ für jeden Teilnehmer.\\
\textbf{Stufe 2} & Regression von $y$ auf $x$ \kursiv{und} $\hat\lambda$, nur für Teilnehmer.\\
\bottomrule
\end{tabular}}
\end{center}

\schritt{Im Beispiel}
\begin{center}
{\small
\renewcommand{\arraystretch}{1.48}
\begin{tabular}{@{}P{3.0cm}P{2.2cm}P{2.4cm}@{}}
\toprule
\textbf{Verfahren} & \textbf{Bildung} & $\hat\lambda$\textbf{-Koeff.}\\
\midrule
KQ & $0{,}048$ & ---\\
Heckman & $\mathbf{0{,}080}$ & $0{,}344$\\
\midrule
wahr & $0{,}080$ & $\rho\sigma = 0{,}350$\\
\bottomrule
\end{tabular}}
\end{center}

\begin{center}
\begin{minipage}{0.88\textwidth}
\kursiv{\textbf{Der Koeffizient von $\hat\lambda$ ist selbst aufschlussreich}: Er schätzt $\rho\sigma$ und sagt damit, ob und wie stark die Auswahl verzerrt. Ist er null, hätten die kleinsten Quadrate genügt. Ein Test auf diesen Koeffizienten ist der übliche Test auf Auswahlverzerrung --- und einer der wenigen Fälle, in denen man eine Endogenität unmittelbar prüfen kann.}
\end{minipage}
\end{center}
\end{merksatz}

\section{Der Bruchfall}

\begin{warnung}
\boxtitelr{Wovon die Korrektur lebt}
Im Beispiel stand die Kinderzahl in der Teilnahmegleichung, nicht in der Lohngleichung. Nun nehme man sie heraus:

\begin{center}
{\small
\renewcommand{\arraystretch}{1.48}
\begin{tabular}{@{}P{3.4cm}P{2.2cm}P{2.6cm}@{}}
\toprule
\textbf{Teilnahmegleichung} & \textbf{Bildung} & $\operatorname{Corr}(x,\hat\lambda)$\\
\midrule
mit Kinderzahl & $0{,}080$ & $-0{,}85$\\
ohne Kinderzahl & $0{,}085$ & $\mathbf{-0{,}997}$\\
\bottomrule
\end{tabular}}
\end{center}

\begin{center}
\begin{minipage}{0.9\textwidth}
\textbf{Ohne eine Größe, die nur die Teilnahme bestimmt, ist $\hat\lambda$ fast eine Funktion von $x$ --- und die Korrektur beruht allein auf der Krümmung von $\lambda$.}
\end{minipage}
\end{center}

\begin{center}
\begin{minipage}{0.88\textwidth}
\kursiv{\textbf{Das ist die Ausschlussbedingung aus Kapitel~2 in neuem Gewand.} Formal ist das Modell auch ohne sie identifiziert, weil $\lambda$ nicht linear ist. Praktisch hängt dann alles an der Normalverteilungsannahme, die diese Krümmung erzeugt --- und an nichts sonst. Heckman selbst hat diese Anwendung seines Verfahrens als unseriös bezeichnet. \textbf{Man braucht ein Instrument für die Teilnahme}, so wie Kapitel~3 eines für den Regressor brauchte.}
\end{minipage}
\end{center}

\schritt{Und die Normalverteilung}
\begin{center}
\begin{minipage}{0.88\textwidth}
\kursiv{Die Korrektur ist bei anderer Störverteilung inkonsistent --- dieselbe Lage wie beim Tobit. Es gibt halbparametrische Fassungen, die $\lambda$ durch eine flexible Funktion des Teilnahmeindex ersetzen; sie brauchen die Ausschlussgröße dann zwingend. \textbf{Je weniger man über die Verteilung annimmt, desto mehr muss das Instrument tragen.}}
\end{minipage}
\end{center}
\end{warnung}

\section{Auswahl ohne Formel}

\begin{merksatz}
\boxtitel{Wo Auswahl sonst noch sitzt}
\begin{center}
{\small
\renewcommand{\arraystretch}{1.52}
\begin{tabular}{@{}P{3.0cm}P{5.6cm}@{}}
\toprule
\textbf{Antwortausfall} & Wer den Fragebogen nicht ausfüllt, unterscheidet sich von denen, die es tun.\\
\textbf{Studienabbruch} & Wer eine Studie verlässt, geht meist wegen des Verlaufs.\\
\textbf{Überleben} & Wer die Nachbeobachtung erreicht, hat die Zwischenzeit überlebt.\\
\textbf{Veröffentlichung} & Studien mit Befund werden gedruckt, Studien ohne nicht.\\
\bottomrule
\end{tabular}}
\end{center}

\begin{center}
\begin{minipage}{0.88\textwidth}
\kursiv{\textbf{In Band 23 hießen diese Fälle Kollider}: Man bedingt auf eine Größe, die von Ursache und Wirkung gemeinsam abhängt. Heckmans Modell gibt dem Kollider eine Gleichung und macht ihn damit rechenbar --- um den Preis der Verteilungsannahme. Band 23, Kapitel~5 hat denselben Schaden als Graph beschrieben; dieses Kapitel beziffert ihn.}
\end{minipage}
\end{center}
\end{merksatz}

\section{Übungsaufgaben}

\begin{aufgabe}
\boxtitel{Aufgabe 6.1 --- Die Richtung der Auswahlverzerrung}
Wann unterschätzt, wann überschätzt die Regression auf die Teilnehmer?

\loesung
\schritt{Schritt 1 --- Der Zusatzterm}
\[ \rho\sigma\,\lambda(z^{\top}\gamma) \quad\text{mit}\quad \lambda \ \text{fallend in } z^{\top}\gamma \]

\schritt{Schritt 2 --- Die Fallunterscheidung}
\begin{center}
{\small
\renewcommand{\arraystretch}{1.5}
\begin{tabular}{@{}P{1.6cm}P{3.2cm}P{3.4cm}@{}}
\toprule
$\rho$ & $x$ \textbf{erhöht Teilnahme} & \textbf{Verzerrung}\\
\midrule
$> 0$ & ja & nach unten\\
$> 0$ & nein & nach oben\\
$< 0$ & ja & nach oben\\
$< 0$ & nein & nach unten\\
\bottomrule
\end{tabular}}
\end{center}

\schritt{Schritt 3 --- Im Beispiel}
Bildung erhöht die Teilnahme, $\rho = 0{,}7 > 0$: Verzerrung nach unten, $0{,}048$ statt $0{,}080$. \checkmark

\schritt{Schritt 4 --- Ein klinischer Fall}
\begin{center}
\begin{minipage}{0.88\textwidth}
\kursiv{Eine Therapiestudie misst den Erfolg nur bei denen, die bis zum Ende bleiben. Wirkt die Therapie ($x$), bleiben mehr; und wer trotz Nebenwirkungen bleibt, hat vermutlich mehr Nutzen ($\rho > 0$). \textbf{Erste Zeile: Die Wirkung wird unterschätzt} --- die Studie ist konservativ. Wird dagegen die Kontrollgruppe häufiger verlassen, weil es ihr schlechter geht, kehrt sich das Vorzeichen um. Ohne die Tabelle kann man das nicht wissen; mit ihr kann man es zumindest begründen.}
\end{minipage}
\end{center}
\end{aufgabe}

\begin{aufgabe}
\boxtitel{Aufgabe 6.2 --- Die Lohnlücke zerlegen}
Erwerbstätige verdienen im Beispiel im Mittel $2{,}37$ (logarithmisch), die Gesamtbevölkerung hätte $2{,}04$. Woraus besteht die Differenz?

\loesung
\schritt{Schritt 1 --- Zwei Quellen}
\[ \mathbb{E}[y\mid d=1] - \mathbb{E}[y] = \underbrace{\beta\bigl(\mathbb{E}[x\mid d=1]-\mathbb{E}[x]\bigr)}_{\text{Zusammensetzung}} + \underbrace{\rho\sigma\,\mathbb{E}[\lambda\mid d=1]}_{\text{Auswahl}} \]

\schritt{Schritt 2 --- Der Auswahlanteil}
\[ 0{,}35\cdot0{,}666 = 0{,}233 \]

\schritt{Schritt 3 --- Der Rest}
\[ 0{,}33 - 0{,}233 = 0{,}10 \quad\text{aus höherer Bildung der Erwerbstätigen} \]

\begin{center}
\begin{minipage}{0.9\textwidth}
\textbf{Siebzig Prozent des Lohnvorsprungs der Erwerbstätigen stammen nicht aus ihren Merkmalen, sondern aus der Auswahl.}
\end{minipage}
\end{center}

\begin{center}
\begin{minipage}{0.88\textwidth}
\kursiv{\textbf{Für die Politik heißt das}: Wer Nichterwerbstätige in Arbeit bringt, darf nicht mit den Löhnen der Erwerbstätigen rechnen. Die Neuankömmlinge werden im Mittel deutlich weniger verdienen --- nicht, weil die Maßnahme versagt, sondern weil die bisherigen Erwerbstätigen eine ausgewählte Gruppe sind. Dieselbe Rechnung gilt für Patienten, die man neu in ein Behandlungsprogramm aufnimmt.}
\end{minipage}
\end{center}
\end{aufgabe}

\begin{aufgabe}
\boxtitel{Aufgabe 6.3 --- Schranken ohne Modell}
Wie weit kann der mittlere Lohn der Gesamtbevölkerung liegen, wenn man über die Nichterwerbstätigen \kursiv{nichts} annimmt?

\loesung
\schritt{Schritt 1 --- Die Zerlegung}
\[ \mathbb{E}[y] = \Pr[d=1]\,\mathbb{E}[y\mid d=1] + \Pr[d=0]\,\mathbb{E}[y\mid d=0] \]
Der letzte Faktor ist unbeobachtet.

\schritt{Schritt 2 --- Die Schranken}
Liegt $y$ sicher in $[y_{\min},y_{\max}]$, so gilt
\[ \Pr[d=1]\,\bar y_1 + \Pr[d=0]\,y_{\min} \ \le\ \mathbb{E}[y] \ \le\ \Pr[d=1]\,\bar y_1 + \Pr[d=0]\,y_{\max} . \]

\schritt{Schritt 3 --- Im Beispiel}
Mit Erwerbsquote $0{,}46$ und plausiblen Grenzen $[1{,}0;\ 3{,}5]$ für den logarithmischen Lohn:
\[ 0{,}46\cdot2{,}37 + 0{,}54\cdot1{,}0 = 1{,}63 \ \le\ \mathbb{E}[y] \ \le\ 0{,}46\cdot2{,}37+0{,}54\cdot3{,}5 = 2{,}98 \]

\begin{center}
\begin{minipage}{0.9\textwidth}
\textbf{Ohne Annahme über die Nichterwerbstätigen weiß man nur: zwischen $1{,}6$ und $3{,}0$. Der wahre Wert $2{,}04$ liegt darin --- und fast alles andere auch.}
\end{minipage}
\end{center}

\schritt{Schritt 4 --- Was das zeigt}
\begin{center}
\begin{minipage}{0.88\textwidth}
\kursiv{\textbf{Das Intervall ist die ehrliche Antwort aus dem Forschungskasten von Kapitel~2.} Heckmans Modell verengt es auf einen Punkt --- aber jede Verengung stammt aus der Annahme, nicht aus den Daten. Die Breite des Intervalls ist das Maß dafür, wie viel das Modell beiträgt: hier fast alles. Wer das weiß, berichtet den Punkt mit der Annahme, nicht ohne sie.}
\end{minipage}
\end{center}
\end{aufgabe}

\begin{geschichte}
\boxtitelg{Ein Aufsatz über Frauenlöhne}
\Person{James Heckman} veröffentlichte das Verfahren 1976 und 1979 --- ursprünglich, um das Arbeitsangebot verheirateter Frauen zu schätzen, deren Löhne man nur bei den Erwerbstätigen kannte.

\begin{center}
\begin{minipage}{0.9\textwidth}
\kursiv{Der Aufsatz von 1979 trägt den Titel \kursiv{Sample Selection Bias as a Specification Error} --- \textbf{Auswahlverzerrung als Spezifikationsfehler}: genau die Einsicht von Satz 6.2, dass die Auswahl wie eine ausgelassene Größe wirkt. Er gehört zu den meistzitierten Arbeiten der Ökonometrie, und seine Anwendung wurde so verbreitet, dass Heckman selbst 1990 vor ihr warnte.}
\end{minipage}
\end{center}

\schritt{Der Nobelpreis}
\begin{center}
\begin{minipage}{0.9\textwidth}
\kursiv{2000 erhielt Heckman den Preis gemeinsam mit \Person{McFadden} --- für die Entwicklung von Methoden zur Analyse ausgewählter Stichproben. \textbf{In seiner Preisrede warnte er vor der Mechanik}: Das Verfahren sei kein Ersatz für ein glaubwürdiges Instrument, sondern verlange eines. Die Warnung wird bis heute regelmäßig überhört.}
\end{minipage}
\end{center}

\schritt{Und die Epidemiologie}
\begin{center}
\begin{minipage}{0.9\textwidth}
\kursiv{Die Medizin kannte das Problem unter anderen Namen --- Berksons Verzerrung von 1946, Überlebensverzerrung, Verlust in der Nachbeobachtung. \textbf{Die beiden Fächer bemerkten erst in den 2000er Jahren, dass sie dasselbe beschrieben}, als die Graphen aus Band 23 eine gemeinsame Sprache lieferten. Seither heißen alle diese Fälle Bedingen auf einen Kollider.}
\end{minipage}
\end{center}
\end{geschichte}

\begin{forschung}
\boxtitelp{Auswahl heute}
\begin{center}
{\small
\renewcommand{\arraystretch}{1.45}
\begin{tabular}{@{}P{3.2cm}P{5.0cm}@{}}
\toprule
\textbf{Schranken} & Manskis Intervalle mit schwachen Zusatzannahmen --- Monotonie, Ordnung.\\
\textbf{Halbparametrisch} & Korrekturen ohne Normalverteilung, mit Instrumenten.\\
\textbf{Zensierte Zielgrößen} & Auswahl in Überlebenszeiten und Panels.\\
\bottomrule
\end{tabular}}
\end{center}

\schritt{Zur ersten Zeile}
\begin{center}
\begin{minipage}{0.88\textwidth}
\kursiv{Nimmt man an, dass die Nichterwerbstätigen im Mittel \kursiv{höchstens} so viel verdient hätten wie die Erwerbstätigen, so halbiert sich das Intervall aus Aufgabe 6.3 --- eine schwache Annahme mit großer Wirkung. \textbf{Die Kunst besteht darin, die schwächste Annahme zu finden, die das Intervall brauchbar macht}, und sie zu benennen. Das ist der Kern von Manskis Programm und die methodische Gegenposition zu Heckmans Punktschätzung.}
\end{minipage}
\end{center}
\end{forschung}

\begin{merksatz}
\boxtitel{Was Sie jetzt können}
\begin{center}
\renewcommand{\arraystretch}{1.5}
\begin{tabular}{@{}P{3.0cm}P{5.6cm}@{}}
\toprule
\textbf{Definition 6.1} & zwei Gleichungen, korrelierte Störungen\\
\textbf{Satz 6.2} & Auswahl als berechenbare ausgelassene Größe\\
\textbf{Zweistufig} & Probit, dann Regression mit $\hat\lambda$\\
\textbf{Der Bruchfall} & ohne Ausschlussgröße trägt nur die Krümmung\\
\textbf{Schranken} & was ohne Annahme bleibt\\
\bottomrule
\end{tabular}
\end{center}

\begin{center}
\begin{minipage}{0.88\textwidth}
\kursiv{Damit endet Teil~II. Alle bisherigen Daten waren Querschnitte --- jeder Versicherte einmal. Teil~III beobachtet dieselben Personen mehrfach und findet, dass sich ein Teil des Problems aus Kapitel~1 dadurch von selbst erledigt.}
\end{minipage}
\end{center}
\end{merksatz}
